\TieudeT{PHONG TỎA VẬT LÍ}
\subsubsection*{PHẦN I. Thí sinh trả lời câu hỏi từ câu 1 đến câu 18. Mỗi câu hỏi thí sinh chỉ chọn một phương án}
\setcounter{ex}{0}
\begin{ex}%book%[1P2Y3-1][Loai1] 
Hệ thức nào sau đây là định luật 2 Newton?
\choice
{$-\vec F = m\vec a$}
{$\vec F = - m\vec a$}
{\True $\vec F = m\vec a$}
{$F=\vec F = ma$}
\loigiai{
Biểu thức của định luật II Niu tơn: $\vec F = m\vec a$.
}
\haidong{8} \end{ex}
\begin{ex}%book%[1P2Y3-1][Loai1] 
Khi nói về một vật chịu tác dụng của lực, phát biểu nào sau đây là đúng?
\choice
{\True Gia tốc của vật luôn cùng chiều với chiều của lực tác dụng}
{Khi ngừng tác dụng lực lên vật, vật này sẽ dừng lại}
{Khi không có lực tác dụng, vật không thể chuyển động}
{Khi có lực tác dụng lên vật, vận tốc của vật tăng}
\loigiai{
Gia tốc của vật luôn cùng chiều với chiều của lực tác dụng.
}
\haidong{8} \end{ex}
\begin{ex}%book%[1P2Y3-1][Loai1]
Dưới tác dụng của một lực vật đang thu gia tốc; nếu lực tác dụng lên vật giảm đi thì độ lớn gia tốc sẽ
\choice
{tăng lên}
{\True giảm đi}
{không đổi}
{bằng $0$}
\loigiai{ 
Khi đó gia tốc của vật giảm đi do a tỉ lệ thuận với F (với m không đổi).
}
\haidong{8} \end{ex}
\begin{ex}%book %[1P2Y3-1][Loai1]
Nếu một vật đang chuyển động có gia tốc mà lực tác dụng lên vật giảm đi thì vật sẽ thu được gia tốc
\choice
{lớn hơn}
{\True nhỏ hơn}
{không thay đổi}
{bằng $0$}
\loigiai{ 
Khi đó gia tốc của vật nhỏ hơn do $F = ma$ (với m không đổi).
}
\haidong{8} \end{ex}
\begin{ex}%book
Phát biểu nào sau đây đúng?
\choice
{Nếu thôi không tác dụng lực vào vật thì vật đang chuyển động sẽ dừng lại}
{\True Nếu chỉ có một lực tác dụng lên vật thì vận tốc của vật bị thay đổi}
{Vật nhất thiết phải chuyển động theo hướng của lực tác dụng}
{Nếu không có lực tác dụng vào vật thì vật không chuyển động được}
\loigiai{\begin{itemize}
\item Chọn B
\end{itemize}

} \haidong{8} \end{ex}
\begin{ex}%book
Một chất điểm đang chuyển động thẳng đều thì chịu tác dụng của một lực không đổi. Phát biểu nào sau đây đúng?
\choice
{Vận tốc của vật sẽ tăng}
{Vận tốc của vật sẽ giảm}
{\True Vận tốc của vật sẽ thay đổi}
{Vận tốc của vật không thay đổi}
\loigiai{\begin{itemize}
\item Chọn C
\end{itemize}

} \haidong{8} \end{ex}
\begin{ex}%book%[1P2Y3-1][Loai1]
Phát biểu nào sau đây là \textbf{sai} khi nói về định luật 2 
Newton? 
\choice
{Với cùng một lực, gia tốc thu được tỉ lệ nghịch với khối lượng của vật}
{Gia tốc mà vật nhận được luôn cùng hướng với lực tác dụng}
{\True Vật luôn chuyển động cùng hướng của lực tác dụng}
{Với cùng một vật, gia tốc thu được tỉ lệ thuận với lực tác dụng}
\loigiai{
Trong trường hợp vật chịu tác dụng của lực ma sát thì vật chuyển động ngược hướng với hướng của lực tác dụng.
}
\haidong{8} \end{ex}
\begin{ex}%book%[1P2B2-1][Loai1]
Đại lượng đặc trưng cho mức quán tính của một vật là
\choice
{trọng lượng}
{\True khối lượng}
{vận tốc}
{lực}
\haidong{8} \end{ex}
\begin{ex}%book
Kết luận nào sau đây chính xác nhất?
\choice
{Vật có khối lượng càng lớn thì rơi càng nhanh}
{Khối lượng riêng của một vật tuỳ thuộc vào khối lượng vật đó}
{\True Vật có khối lượng càng lớn càng khó thay đổi vận tốc}
{Để đo khối lượng người ta dùng lực kế}
\loigiai{
Khối lượng liên quan tới quán tính, khối lượng càng lớn thì quán tính càng lớn càng khó thay đổi vận tốc ban đầu.
} \haidong{8} \end{ex}
\begin{ex}%book
Lực $\vec{F}$ truyền cho vật khối lượng ${m}_1$ gia tốc ${a}_1$, truyền cho vật khối lượng ${m}_2$ gia tốc ${a}_2$. Lực $\vec{F}$ sẽ truyền cho vật có khối lượng $m=m_1+m_2$ gia tốc
\choice
{$\dfrac{a_1+a_2}{2}$}
{$\dfrac{a_1+a_2}{a_1 a_2}$}
{\True $\dfrac{a_1 a_2}{a_1+a_2}$}
{$a_1+a_2$}
\loigiai{\begin{itemize}
\item 	Chọn C
\end{itemize}

} \haidong{8} \end{ex}
\begin{ex}%book
Một ô tô có khối lượng $1600\ kg$ đang chuyển động thì bị hãm phanh với lực hãm có độ lớn bằng $600\ N$. Vectơ gia tốc mà lực này gây ra cho xe có độ lớn và đặc điểm nào sau đây?
\choice
{$0,375\ m/s^2$, cùng với hướng chuyển động}
{\True $0,375\ m/s^2$, ngược với hướng chuyển động}
{$\dfrac{8}{3}\ m/s^2$, ngược với hướng chuyển động}
{$\dfrac{8}{3}\ m/s^2$, cùng với hướng chuyển động}
\haidong{5}\loigiai{\begin{itemize}
\item Chọn chiều dương là chiều chuyển động.
\item Từ: $a=\dfrac{F_{h l}}{m}=\dfrac{-600}{1600}=-0,375\ m/s^2$
\item Chọn B
\end{itemize}

} \haidong{8} \end{ex}
\begin{ex}%book%[1P2K3-2][Loai1]
Một chiếc xe khối lượng $m=100\ kg$ đang chạy với vận tốc 30,6 km/h thì hãm phanh. Biết lực hãm là $250\ N$. Quãng đường xe còn chạy thêm trước khi dừng hẳn là
\choice
{\True $14,45\ m$}
{$24,45\ m$}
{$44,35\ m$}
{$64,65\ m$}
\haidong{5}\loigiai{
Lực tác dụng lên xe là lực hãm. Theo định luật II Niutơn ta có: \bth{\vec{F}=m\vec{a}}\hfill (1)\\ Chiếu (1) lên hướng chuyển động ta có: \bth{-F=ma}
\begin{itemize}
\item Gia tốc chuyển động: \bs{a}{-\dfrac{F}{m}=-2,5\ m/s^2}
\item Khi xe bắt đầu hãm phanh: \bs{v_0}{30,6\ km/h=8,5\ m/s}
\item Khi xe dừng: \bs{v}{0}
\end{itemize}
Quãng đường xe chạy thêm: \bs{s}{\dfrac{v^2-v_0^2}{2a}=14,45\ m}.
}
%<MyLT>
\haidong{8} \end{ex}
\begin{ex}%book
Một ô tô đang chạy với tốc độ $60\ km/h$ thì người lái xe hãm phanh, xe đi tiếp được quãng đường $50\ m$ thì dừng lại, lực hãm trong hai trường hợp bằng nhau. Nếu ô tô chạy với tốc độ $120\ km/h$ thì quãng đường đi được từ lúc hãm phanh đến khi dừng lại là
\choice
{$70,7\ m$}
{$141\ m$}
{$100\ m$}
{\True $200\ m$}
\haidong{9}\loigiai{\begin{itemize}
\item Chọn chiều dương là chiều chuyển động. Vì lực hãm trong hai trường hợp bằng nhau nên độ lớn gia tốc bằng nhau và bằng ${a}$.
\item 	$
0^2-v_0^2=-2 \text {as} \Rightarrow \dfrac{v_{02}^2}{v_{01}^2}=\dfrac{s_2}{s_1} \Rightarrow s_2=s_1\left(\dfrac{v_{02}}{v_{01}}\right)^2=50.2^2=200\ m
$
\item Chọn D
\end{itemize}

} \haidong{8} \end{ex}
\begin{ex}%book
Một quả bóng có khối lượng $500\ g$ đang nằm trên mặt đất thì bị đá bằng một lực $200\ N$. Nếu thời gian quả bóng tiếp xúc với bàn chân là $0,02\ s$ thì bóng sẽ bay đi với tốc độ bằng
\choice
{$2\ m/s$}
{$0,8\ m/s$}
{$0,008\ m/s$}
{\True $8\ m/s$}
\haidong{5}\loigiai{\begin{itemize}
\item Theo định luật II Niu - tơn: $a=\dfrac{F}{m}=\dfrac{200}{0,5}=400\ m/s^2$
\item Thời gian bóng tiếp xúc với bàn chân là thời gian bóng được truyền gia tốc. Vậy vận tốc của bóng khi bay đi là: $v=v_0+a t=0+400.0,02=8\ m/s$
\item Chọn D
\end{itemize}

} \haidong{8} \end{ex}
\begin{ex}%book
Một lực không đổi tác dụng vào một vật có khối lượng $5,0\ kg$ làm tốc độ của nó tăng dần từ $2,0\ m/s$ đến $8,0\ m/s$ trong $3,0\ s$. Độ lớn lực tác dụng vào vật là
\choice
{$15\ N$}
{\True $10\ N$}
{$1,0\ N$}
{$5,0\ N$}
\haidong{5}\loigiai{\begin{itemize}
\item Chọn chiều dương là chiều chuyển động.
\item Từ: $\left\{\begin{array}{l}a=\dfrac{v_s-v_t}{\Delta t}=\dfrac{8-2}{3}=2\ m/s^2 \\ F_{h l}=m a=5.2=10\ N\end{array}\right.$
\item Chọn B
\end{itemize}

} \haidong{8} \end{ex}
\begin{ex}%book
Phải tác dụng một lực $50\ N$ theo hướng chuyển động vào một xe chở hàng có khối lượng $400\ kg$ trong thời gian bao nhiêu để tăng tốc độ của nó từ $10\ m/s$ lên đến $12\ m/s$?
\choice
{$10\ s$}
{$20\ s$}
{\True $16\ s$}
{$40\ s$}
\haidong{5}\loigiai{\begin{itemize}
\item Chọn chiều dương là chiều chuyển động.
\item Từ: $a=\dfrac{F}{m}=\dfrac{50}{400}=0,125\ m/s^2 \Rightarrow t=\dfrac{v_s-v_t}{a}=\dfrac{12-10}{0,125}=16\ s$
\item Chọn C
\end{itemize}

} \haidong{8} \end{ex}
\begin{ex}%book%[1P2K3-2][Loai1] 
Một ôtô có khối lượng $3$ tấn, đang chạy với vận tốc $v_0$ thì hãm phanh, xe đi thêm quãng đường $15\ m$ trong $3\ s$ thì dừng hẳn. Bỏ qua các lực cản bên ngoài. Vận tốc $v_0$ và độ lớn lực hãm phanh lần lượt là
\choice
{$10\ m/s$ và $2.10^4\ N$}
{$15\ m/s$ và $10^4\ N$}
{\True $10\ m/s$ và $10^4\ N$}
{$15\ m/s$ và $2.10^4\ N$}
\haidong{5}\loigiai{
\begin{itemize}
\item Ta có: $\begin{cases}
v=v_0+at\\
S=v_0t+0,5at^2\\
\end{cases}\Rightarrow \begin{cases}
0=v_0+3a\\
15=3v_0+4,5a\\
\end{cases}\Rightarrow \begin{cases}
v_0=10\ m/s \\
a=-\dfrac{10}{3}\ m/s^2 \\
\end{cases}$
\item Chọn chiều dương là chiều chuyển động của xe
\begin{itemize}
\item Biểu thức định luật II Niu-tơn: $\vec{F}_{h}=m\vec{a}$
\item Chiếu lên chiều dương ta có: $-{F}_{h}=ma \Leftrightarrow -F_h={3.10}^{3}\left(-\dfrac{10}{3}\right)\Rightarrow F_h={10}^{4}\ N$
\item Vậy độ lớn của lực hãm là $F_h={10}^{4}\ N$. 
\end{itemize}
\end{itemize}
}
\haidong{8} \end{ex}

\begin{ex}%book
Một vật có khối lượng $m$, chịu tác dụng đồng thời của hai lực có độ lớn lần lượt là $F_1=4\ N$ và $F_2=6\ N$. Gia tốc của vật có độ lớn nhỏ nhất trong trường hợp
\choice
{$\vec{F}_1$ vuông góc với $\vec{F}_2$}
{$\vec{F}_1$ cùng chiều với $\vec{F}_2$}
{$\vec{F}_1$ hợp với $\vec{F}_2$ một góc $60^{\circ}$}
{\True $\vec{F}_1$ ngược chiều với $\vec{F}_2$}
\haidong{5}\loigiai{
\begin{itemize}
\item $|{a}|=\dfrac{|{F}|}{{m}}$, do ${m}$ không đổi nên $|{a}|$ nhỏ nhất khi $|{F}|$ nhỏ nhất $\Leftrightarrow$ hai lực cùng phương, ngược chiều nhau.
\item Chọn D
\end{itemize}	

} \haidong{8} \end{ex}
\subsubsection*{PHẦN 2. Thí sinh trả lời từ câu 1 đến câu 4. Trong mỗi ý a), b), c), d) ở mỗi câu, thí sinh chọn đúng hoặc sai.}
\setcounter{ex}{0}
\begin{ex}%DoanhNamDinh
\immini[thm]{Cho đồ thị biểu diễn mối liên hệ giữa các lực tác dụng lên một vật và gia tốc do lực gây ra tương ưng như hình vẽ.}{
\includegraphics[width=4cm]{images/2025_06_04_af03f1d88c157fe9bf98g-138}}
\choiceTF[t]
{\True Khi lực có độ lớn $1\ N$ tác dụng lên vật thì gia tốc của vật là $2 \ m/s^2$}
{\True Đồ thị là đường thẳng xiên góc đi qua gốc tọa độ nên độ lớn của gia tốc $a$ tỉ lệ thuận với độ lớn của lực $F$ tác dụng lên vật}
{Khối lượng của vật bằng $0,5\ kg$ nên góc tạo bởi đồ thị và trục hoành bằng $0,5^\circ$}
{Nếu tác dụng đồng thời lên vật hai lực có độ lớn $1\ N$ và $2\ N$ thì gia tốc của vật là $6 \ m/s^2$}
\loigiai{
\begin{itemchoice}
\itemch Dựa vào đồ thị ta có $F=1 \ N$ thì $a=2 \ m/s^2$
$\Rightarrow$ Đúng
\itemch Dựa vào nội dung định luật 2 Newton.
$\Rightarrow$ Đúng
\itemch Khối lượng của vật $m=F / a=0,5 \ kg=\tan \alpha$.
$\Rightarrow$ Sai
\itemch Vì chưa cho biết hai lực có hướng như thế nào nên chưa thể biết độ lớn của hợp lực.
$\Rightarrow$ Sai
\end{itemchoice}
}
\end{ex}

\begin{ex}%DoanhNamDinh
Một vật có khối lượng $20\ kg$ đang đứng yên thì chịu thêm tác dụng đồng thời của hai lực $\vec{F}_1$ và $\vec{F}_2$ có cùng giá, ngược chiều, độ lớn lần lượt là $40\ N$ và $30\ N$, làm vật chuyển động thẳng theo hướng của hợp lực.
\choiceTF[t]
{Vật ở trạng thái cân bằng}
{\True Gia tốc của vật cùng hướng với hướng của lực $\vec{F}_1$}
{Gia tốc của vật có độ lớn là $3,5 \ m/s^2$}
{\True Tốc độ của vật sau $5\ s$ kể từ lúc bắt đầu chuyển động là $2,5 \ m/s$}
\loigiai{
\begin{itemchoice}
\itemch Do $F_1>F_2$ nên hợp lực tác dụng lên vật không bằng o.
$\Rightarrow$ Sai
\itemch Do $F_1>F_2$ nên hợp lực tác dụng lên vật cùng hướng với lực $\vec{F}_1$, gia tốc luôn cùng hướng với hợp lực. $\Rightarrow$ Đúng
\itemch Gia tốc của vật có độ lớn là $(4 O-3 O): 2 O=0,5 \ m/s^2$.
$\Rightarrow$ Sai
\itemch Tốc độ của vật sau $5\ s$ là $v=a. t=0,5 \cdot 5=2,5 \ m/s$.
$\Rightarrow$ Đúng
\end{itemchoice}
}
\end{ex}

\begin{ex}%DoanhNamDinh
\impicinpar{Một xe có khối lượng $1,5$ tấn đang chuyển động thẳng đều với tốc độ $72 \ km / h$ thì tắt máy, chuyển động thẳng chậm dần đều trên đường nằm ngang dưới tác dụng của lực cản. Biết lực cản tác dụng lên xe không đổi, luôn ngược hướng chuyển động của xe và có độ lớn $F_c=3000 \ N$.}{
\includegraphics[width=4cm]{images/2025_06_04_af03f1d88c157fe9bf98g-135(1)}}
\choiceTF[t]
{Sau khi xe tắt máy, vectơ gia tốc của xe luôn cùng hướng với vectơ vận tốc}
{Gia tốc của xe có độ lớn $3 \ m/s^2$}
{\True Quãng đường xe đi được từ lúc tắt máy cho đến khi dừng hẳn là $100\ m$}
{Thời gian kể từ lúc xe tắt máy đến lúc dừng hẳn là $20\ s$}
\loigiai{
\begin{itemchoice}
\itemch Xe chuyển động thẳng chậm dần đều nên gia tốc của xe luôn ngược hướng với vectơ vận tốc.
$\Rightarrow$ Sai
\itemch Gia tốc của xe có độ lớn $|a|=\left|\dfrac{F}{m}\right|=3000: 1500=2 \ m/s^2$.
$\Rightarrow$ Đúng
\itemch Quãng đường xe đi được cho đến khi dừng lại $s=\dfrac{v^2-v_{O}^2}{2 a}=100 \ m$
$\Rightarrow$ Đúng
\itemch Thời gian kể từ lúc xe tắt máy đến lúc dừng lại là $v=v_0+a t \leftrightarrow 0=20-2 t \rightarrow t=10\ s$
$\Rightarrow$ Sai
\end{itemchoice}
}
\end{ex}
\begin{ex}%DoanhNamDinh
\immini[thm]{Một vật có khối lượng $m=4\ kg$ chuyển động theo một đường thẳng và có đồ thị vận tốc theo thời gian được biểu diễn như hình vẽ.}{\begin{tikzpicture}[line cap=round, line join=round, scale=1, font=\footnotesize, >=stealth,draw=Mapcolor]
%[draw=Mapcolor]
\path
(0,0)coordinate(O)
(5,0)coordinate(x)
(0,3.5)coordinate(y)
;
\draw  [dashed,help  lines, xstep=0.5cm,ystep=0.5cm]  (0,0)  grid  (4.5,3);
\draw [->](O)--(x);
\draw [->](O)--(y);

\foreach  \x  in  {1,2,...,4}
{\pgfmathtruncatemacro{\label}{2*\x};
\node[below]
(\label) at (\x,0) {\label};
\draw[fill=white](\x,0)circle(0.5pt);}

\foreach  \y  in  {0,1,2,3}
{\pgfmathtruncatemacro{\label}{2*\y};
\node[left]
(\label) at (0,\y) {\label};
\draw[fill=white](0,\y)circle(0.5pt);}
\node at (0,4) {$v\ (m/s)$};
\node at (5.5,0) {$t\ (s)$};
\draw[very thick,Mapcolor] (0,0) -- (1,2.5) -- (3.5,2.5) -- (4,0);
\end{tikzpicture}}
\choiceTF[t]
{\True Từ thời điểm $t_1=2 \ s$ đến $t_2=7 \ s$, hợp lực tác dụng lên vật bằng $0$}
{Từ thời điểm $t_2=7 \ s$ đến $t_3=8 \ s$, vật chuyển động thẳng chậm dần đều nên hợp lực tác dụng lên vật ngược hướng với gia tốc của vật}
{Độ lớn hợp lực tác dụng lên vật tại thời điểm $t=1 \ s$ là $24\ N$}
{\True Quãng đường vật đi được trong 8 giây đầu là $32,5 \ m$}
\loigiai{
\begin{itemchoice}
\itemch Từ thời điểm $t_1=2 \ s$ đến $t_2=7 \ s$, vật chuyển động thẳng đều, $a=o \rightarrow F=o$
$\Rightarrow$ Đúng
\itemch Hợp lực tác dụng lên vật luôn cùng hướng với gia tốc của vật.
$\Rightarrow$ Sai
\itemch Trong hai giây đầu, gia tốc của vật là: $a=\dfrac{5-o}{2}=2,5 \ m/s^2$.
Độ lớn hợp lực tác dụng lên vật: $F=m a=10 \ N$.
$\Rightarrow$ Sai
\itemch Quãng đường vật đi được bằng diện tích của hình thang giới hạn bởi hệ trục và đồ thị
$s=\dfrac{(8+5) \cdot 5}{2}=32,5 \ m$
$\Rightarrow$ Đúng
\end{itemchoice}
}
\end{ex}














\subsubsection*{PHẦN 3. Thí sinh trả lời từ câu 1 đến câu 6.}
\setcounter{ex}{0} 
\begin{ex}
Một vật khối lượng $2\ kg$ đặt trên mặt phẳng nằm ngang. Khi tác dụng một lực có độ lớn là $1\ N$ theo phương ngang vật bắt đầu trượt trên mặt phẳng nằm ngang. Bỏ qua mọi ma sát và lực cản. Vận tốc của vật sau $4\ s$ có độ lớn là bao nhiêu $m/s$?
\shortans[oly]{2}
\loigiai{}
\end{ex}
\begin{ex}%DoanhNamDinh
Một hợp lực có độ lớn $10\ N$ tác dụng vào một vật có khối lượng $2\ kg$ đang đứng yên làm cho vật chuyển động theo hướng của lực. Quãng đường vật đi được trong giây thứ $2$ là bao nhiêu mét?
\shortans[oly]{7,5}
\loigiai{}
\end{ex}
\begin{ex}%book
Lực có độ lớn $F$ tác dụng lên một vật trong khoảng thời gian $2\ s$ làm tốc độ của nó thay đổi từ $0,8\ m/s$ đến $1\ m/s$. Biết lực đó có độ lớn không đổi và có phương luôn cùng phương với chuyển động. Nếu lực đó tác dụng lên vật trong khoảng thời gian $1,1\ s$ thì tốc độ của vật thay đổi một lượng là bao nhiêu $m/s$?
\shortans[oly]{0,11}
\loigiai{
\begin{itemize}
\item Chọn chiều dương là chiều chuyển động của vật

$
F=m a=m \dfrac{\Delta v_1}{\Delta t_1}=m \dfrac{\Delta v_2}{\Delta t_2} \Rightarrow \Delta v_2=\Delta v_1 \dfrac{\Delta t_2}{\Delta t_1} \Rightarrow \Delta v_2=(1-0,8) \dfrac{1,1}{2}=0,11\ m/s
$
\item Chọn B
\end{itemize}

} \haidong{8} \end{ex}

\begin{ex}%book
Một ôtô có khối lượng $500\ kg$ đang chuyển động thẳng đều thì tắt máy, hãm phanh chuyển động chậm dần đều trong $2\ s$ cuối cùng đi được $1,8\ m$. Độ lớn lực hãm là bao nhiêu niutơn?
\shortans[oly]{450}
\loigiai{\begin{itemize}
\item Gọi ${v}_0$ là vận tốc ban đầu của quãng đường đi $2\ s$ cuối.
\item Ta có:
$
\begin{aligned}
& {v}^2-{v}_0^2=2 . a . s \Rightarrow-{v}_0^2=2 . a . s=3,6 {a}\ (1) \\
& {a}=\dfrac{{v}-{v}_0}{\Delta {t}} \Rightarrow-{v}_0=\text {at}
\end{aligned}
$
\item Từ (1) và (2) ta có. ${a}=-0,9\ m/s^2$
$\Rightarrow {F}={m} . {a}=-450\ N$. Dấu "-" chứng tỏ lực ngược chiều chuyển động (lực hãm).
\item Chọn B
\end{itemize}

} \haidong{8} \end{ex}
\begin{ex}%book
Một vật đang chuyển động dọc theo chiều dương của trục $Ox$ thì một lực không đổi có phương song song với trục $Ox$, tác dụng vào vật trong khoảng thời gian $0,6\ s$ làm vận tốc của nó thay đổi từ $8\ cm/s$ đến $5\ cm/s$. Tiếp đó, tăng độ lớn của lực lên gấp đôi trong khoảng thời gian $2,2\ s$ nhưng vẫn giữ nguyên hướng của lực, sau khi kết thúc khoảng thời gian này thì vận tốc vật đạt được có giá trị là
\shortans[oly]{-17}
\haidong{5}\loigiai{\begin{itemize}
\item Trong giai đoạn đầu: $a_1=\dfrac{v_s-v_t}{\Delta t}=\dfrac{5-8}{0,6}=-5\ cm/s^2$
\item Trong giai đoạn sau: $a_2=2 a_1=-10\ cm/s^2 \Rightarrow v_x=v_t+a_2 t=5+(-10) . 2,2=-17{\ cm} / {s}$
\item 	Chọn A
\end{itemize}		
} \haidong{8} \end{ex}
\begin{ex}%Tailieuchuan
Một vật có khối lượng $25\ kg$ đang chuyển động với vận tốc $5\ m/s$ trên mặt sàn nằm ngang. Vật chịu tác dụng của lực kéo theo phương ngang có độ lớn $32\ N$, lực cản không đổi ngược chiều chuyển động có độ lớn $12\ N$. Tốc độ trung bình của vật trong khoảng thời gian từ giây thứ $5$ đến hết giây thứ $6$ là bao nhiêu $m/s$?
\shortans[oly]{9,4}
\loigiai{}
\end{ex}







\newpage

\TieudeT{PHONG TỎA VẬT LÍ}
\subsubsection*{PHẦN I. Thí sinh trả lời câu hỏi từ câu 1 đến câu 18. Mỗi câu hỏi thí sinh chỉ chọn một phương án}
\setcounter{ex}{0}
\begin{ex}%book%[1P2Y3-1][Loai1] 
Định luật 2 Newton cho biết
\choice
{mối liên hệ giữa khối lượng và vận tốc của vật}
{lực là nguyên nhân gây ra chuyển động}
{mối liên hệ giữa lực và vận tốc của vật}
{\True lực là nguyên nhân làm xuất hiện gia tốc của vật}
\loigiai{
Định luật II cho biết lực là nguyên nhân làm xuất hiện gia tốc của vật.
}
\haidong{8} \end{ex}
\begin{ex}%book%[1P2B3-2][Loai1]
Một vật đang chuyển động dưới tác dụng của lực $F_1$ với gia tốc $a_1$. Nếu tăng lực lên thành $F_2 = 2F_1$ thì gia tốc $a_2$ của vật bằng
\choice
{\True $2a_1$}
{$4a_1$}
{$a_1$}
{$0,5a_1$}
\loigiai{
Ta có: $a_2 = \dfrac{F_2}{m} = \dfrac{2F_1}{\dfrac{F_1}{a_1}} = 2a_1$.
}
\haidong{8} \end{ex}
\begin{ex}%book%[1P2Y3-1][Loai1]
Trong chuyển động thẳng chậm dần đều thì hợp lực tác dụng vào vật
\choice
{cùng chiều với chuyển động}
{cùng chiều với chuyển động và có độ lớn không đổi}
{ngược chiều với chuyển động và có độ lớn nhỏ dần}
{\True ngược chiều với chuyển động và có độ lớn không đổi}
\loigiai{ 
Trong chuyển động chậm dần đều thì gia tốc ngược chiều chuyển động và có độ lớn không đổi nên lực cũng ngược chiều chuyển động và có độ lớn không đổi.
}
\haidong{8} \end{ex}
\begin{ex}%book
Phát biểu nào sau đây là \textbf{sai}?
\choice
{\True Vật chuyển động thẳng nhanh dần đều khi lực tác dụng lên vật có độ lớn tăng dần}
{Lực tác dụng vào một vật càng lớn thì độ lớn gia tốc của vật càng lớn}
{Dưới tác dụng của cùng một lực, vật nào có khối lượng càng lớn thì độ lớn gia tốc của vật càng nhỏ}
{Vectơ gia tốc của một vật luôn cùng hướng với lực gây ra gia tốc đó}
\loigiai{\begin{itemize}
\item Trong chuyển động thẳng biến đổi đều, gia tốc không đổi nên lực tác dụng không đổi.
\item Chọn A
\end{itemize}

} \haidong{8} \end{ex}
\begin{ex}%book
Trong trường hợp nào dưới đây, vật chuyển động theo hướng của hợp lực tác dụng vào vật?
\choice
{Vật chuyển động tròn đều}
{\True Vật chuyển động thẳng nhanh dần đều}
{Vật chuyển động thẳng chậm dần đều}
{Vật chuyển động thẳng đều}
\loigiai{\begin{itemize}
\item $\vec{{a}}$ luôn cùng hướng với $\vec{F}$; chuyển động thẳng nhanh dần đều $\vec{{a}} ; \vec{{v}}$ cùng hướng; chuyển động thẳng chậm dần đều $\vec{{a}} ;\ \vec{{v}}$ ngược hướng.
\item Chọn B
\end{itemize}


} \haidong{8} \end{ex}
\begin{ex}%book %[1P2Y3-1][Loai1]
Phát biểu nào sau đây là \textbf{sai}?
\choice
{Gia tốc của một vật luôn cùng chiều với hợp lực tác dụng lên vật}
{\True Chiều của vectơ gia tốc chỉ chiều chuyển động của vật}
{Gia tốc của vật càng lớn thì vận tốc biến đổi càng nhanh}
{Độ lớn của gia tốc tỉ lệ thuận với lực tác dụng lên vật và li lệ nghịch với khối lượng của vật}
\loigiai{ 
Chiều chuyển động của vật là chiều của vận tốc chứ không phải là chiều của hợp lực (chiều của gia tốc).
}
\haidong{8} \end{ex}
\begin{ex}%book
Khối lượng là đại lượng đặc trưng cho
\choice
{trọng lượng của vật}
{tác dụng làm quay của lực quanh một trục}
{thể tích của vật}
{\True mức quán tính của vật}
\loigiai{
Khối lượng đặc trưng cho mức quán tính của vật.
} \haidong{8} \end{ex}
\begin{ex}%book %[1P2Y2-1][Loai1]
Phát biểu nào sau đây là \textbf{sai} khi nói về tính chất của khối lượng?
\choice
{Khối lượng là đại lượng bất biến đối với mỗi vật}
{Khối lượng có tính chất cộng được}
{\True Khối lượng của vật càng lớn thì mức quán tính của vật càng nhỏ và ngược lại}
{Trong hệ đơn vị SI đơn vị khối lượng là kilôgam (kg)}
\loigiai{ 
Khối lượng đặc trưng cho mức quán tính của vật do đó khối lượng càng lớn thì tính quán tính càng lớn (tính ì) $\Rightarrow$ Phát biểu Khối lượng của vật càng lớn thì mức quán tính cúa vật càng nhỏ và ngược lại là \textbf{sai}.
}
\haidong{8} \end{ex}
\begin{ex}%book %[1P2B3-1][Loai1]
\immini[thm]{Hình vẽ sau là vẽ các lực tác dụng lên máy bay lên thẳng, trong đó $\vec v$ là vectơ vận tốc của máy bay tại thời điểm đang xét, $\vec F$ lực kéo của động cơ, $\vec R$ lực cản, $\vec P$ trọng lực, $\vec Q$ lực nâng máy bay. Khi máy bay tăng tốc thì các cặp lực $\vec P$ và $\vec Q$, $\vec F$ và $\vec R$ phải thỏa mãn điều kiện nào sau đây?
\choice
{$P = Q,\ F = R$}
{\True $P = Q,\ F> R$}
{$P> Q,\ F = R$}
{$P < Q,\ F = R$}}{		\begin{tikzpicture}[scale=1,thick,>=stealth']
\node[inner sep=0pt,yshift=15pt,xshift=-20pt] (russell) at (0,0){\includegraphics[width=8cm]{images/helicopter1.pdf}};
\draw[very thick,red,->](0,0)--++(-90:2)node[black,right]{$\vec{P}$};
\draw[very thick,red,->](0,0)--++(90:2)node[black,left]{$\vec{Q}$};
\draw[very thick,red,->](0,0)--++(0:3)node[black,above]{$\vec{F}$};
\draw[very thick,red,->](0,0)--++(180:3)node[black,below]{$\vec{R}$};
\draw[fill=white](0,0) circle(1.5pt);
\draw[very thick,blue,->](1.7,-0.2)--++(0:1.5)node[black,below]{$\vec{v}$};
\draw[fill=white](1.7,-0.2) circle(1.5pt);
\end{tikzpicture}}
\loigiai{ 
Do máy bay tăng tốc hướng lên phía trước $\Rightarrow$ $P = Q,\ F> R$.
}
\haidong{8} 
\end{ex}
\begin{ex}%book%[1P2Y3-2][Loai1]
Tác dụng một lực theo phương ngang vào vật có khối lượng $2\ kg$, đang đứng yên thì vật này chuyển động nhanh dần đều với gia tốc $0,5\ m/s^2$. Độ lớn của lực này là
\choice
{$2\ N$}
{$20\ N$}
{$10\ N$}
{\True $1\ N$}
\haidong{5}\loigiai{
Độ lớn của lực là: $F = ma = 2.0,5 = 1\ N$.
}
\haidong{8} \end{ex}
\begin{ex}%book
Một lực có độ lớn $1,0\ N$ tác dụng vào một vật có khối lượng $2,0\ kg$ lúc đầu đứng yên, trong khoảng thời gian $2,0\ s$. Quãng đường mà vật đi được trong khoảng thời gian đó là
\choice
{$2,0\ m$}
{$4,0\ m$}
{$0,5\ m$}
{\True $1,0\ m$}
\haidong{5}\loigiai{\begin{itemize}
\item 	Chọn chiều dương là chiều chuyển động.

$
a=\dfrac{F}{m}=\dfrac{1}{2}=0,5\ m/s^2 \Rightarrow s=0,5 a t^2=0,5 . 0,5 . 2^2=1\ m
$
\item Chọn D
\end{itemize}


} \haidong{8} \end{ex}
\begin{ex}%book
Một vật khối lượng $2\ kg$ đang chuyển động với vận tốc $18\ km/h$ thì bắt đầu chịu tác dụng của lực $4\ N$ theo chiều chuyển động. Đoạn đường vật đi được trong $10\ s$ đầu tiên là
\choice
{$120\ m$}
{$160\ m$}
{$175\ m$}
{\True $150\ m$}
\haidong{5}\loigiai{\begin{itemize}
\item Ta có ${v}=18\ km/h=5\ m/s$.
\item Định luật II: $F=m a \Rightarrow a=\dfrac{F}{m}=2\ m/s^2$
\item Quãng đường vật đi được trong $10\ s$ là: $s=v_0 t+\dfrac{1}{2} a t^2=5.10+\dfrac{1}{2} . 2.100=150\ m$
\item Chọn D
\end{itemize}

} \haidong{8} \end{ex}
\begin{ex}%book
Một quả bóng ${m}=400\ g$ đang nằm yên trên mặt đất. Một cầu thủ đá bóng với lực $300\ N$. Thời gian chân tác dụng vào quả bóng là $0,015\ s$; bỏ qua ma sát. Tốc độ của quả bóng lúc bay đi là
\choice
{$22,5\ m/s$}
{$11,25.10^{-3}\ m/s$}
{\True $11,25\ m/s$}
{$11250\ m/s$}
\haidong{5}\loigiai{\begin{itemize}
\item Ta có $: {a}=\dfrac{F}{{m}}=750\ m/s^2 ; {v}={v}_0+{at}=11,25\ m/s$
\item Chọn C
\end{itemize}


} \haidong{8} \end{ex}
\begin{ex}%book
Một ô tô có khối lượng $1500\ kg$ khi khởi hành được tăng tốc bởi một lực $2000\ N$ trong $18$ giây đầu tiên. Tốc độ của xe đạt được ở cuối khoảng thời gian đó là
\choice
{$45\ m/s$}
{$40\ m/s$}
{\True $24\ m/s$}
{$20\ m/s$}
\haidong{5}\loigiai{\begin{itemize}
\item 	Chọn C
\end{itemize}

} \haidong{8} \end{ex}
\begin{ex}%book
Phải tác dụng một lực $100\ N$ hướng theo chiều chuyển động vào một xe chở hàng đang chuyển động thẳng đều có khối lượng $200\ kg$ trong thời gian bao nhiêu để tăng tốc độ của nó từ $10\ m/s$ lên đến $12\ m/s$ ?
\choice
{$16\ s$}
{$8\ s$}
{$10\ s$}
{\True 4 s}
\haidong{5}\loigiai{\begin{itemize}
\item Chọn chiều dương là chiều chuyển động.
\item Từ ${a}=\dfrac{F}{{m}}=\dfrac{100}{200}=0,5\ m/s^2 \Rightarrow {t}=\dfrac{{v}_{{s}}-{v}_{{t}}}{{a}}=\dfrac{12-10}{0,5}=4{~s}$
\item Chọn D
\end{itemize}

} \haidong{8} \end{ex}
\begin{ex}%book
Một vật có khối lượng $3\ kg$ đang chuyển động thẳng đều với tốc độ $2\ m/s$ thì chịu tác dụng của một lực $9\ N$ cùng hướng với hướng chuyển động. Vật sẽ chuyển động $10\ m$ tiếp theo trong thời gian là
\choice
{$10\ s$}
{$4\ s$}
{$1,6\ s$}
{\True $2\ s$}
\haidong{5}\loigiai{\begin{itemize}
\item Chọn chiều dương là chiều chuyển động.
\item $\left\{\begin{array}{l}a=\dfrac{F}{m}=\dfrac{9}{3}=3\ m/s^2 \\ s=v_0 t+0,5 a t^2 \Rightarrow 10=2 t+1,5 t^2 \Rightarrow t=2\ s\end{array}\right.$
\item Chọn D
\end{itemize} 

} \haidong{8} \end{ex}
\begin{ex}%book%[1P2K3-2][Loai1]
Dưới tác dụng của lực $F$ nằm ngang, xe lăn chuyển động không vận tốc đầu, đi được quãng đường $2,5\ m$ trong thời gian $t$. Nếu đặt thêm vật khối lượng $250\ g$ lên xe thì xe chỉ đi được quãng đường $2\ m$ trong thời gian $t$. Bỏ qua ma sát. Khối lượng của xe là
\choice
{$0,5\ kg$}
{$1,5\ kg$}
{\True $1\ kg$}
{$2\ kg$}
\haidong{5}\loigiai{
\begin{itemize}
\item Gọi khối lượng của xe và của vật đặt thêm lần lượt là m, m’
\item Gia tốc của xe tương ứng với mỗi trường hợp là a, a’
\item Theo định luật II, Niu tơn: $F = ma = (m + m’)a’$
\item Mặt khác, quãng đường xe đi trong 2 trường hợp là: 
\begin{align*}
s = \dfrac{1}{2}at^2,\ s' = \dfrac{1}{2}a't^2 \Rightarrow \dfrac{{m + m'}}{m} = \dfrac{a}{{a'}} = \dfrac{s}{{s'}} \Rightarrow \dfrac{{s'}}{{s - s'}}m' = 22,5-2.0,25=1\dfrac{2}{{2,5 - 2}}.0,25 = 1\ kg.
\end{align*}
\end{itemize}
}

\haidong{8} \end{ex}
\begin{ex}%book
Một vật có khối lượng $m=2\ kg$ đặt trên bàn nhẵn nằm ngang chịu tác dụng của hai lực là $F_1=6\ N$ và $F_2=4\ N$ ngược chiều nhau theo phương ngang, bỏ qua ma sát. Gia tốc của vật thu được là
\choice
{$2\ m/s^2$ hướng theo $F_1$}
{$2\ m/s^2$ hướng theo $F_2$}
{\True $1\ m/s^2$ hướng theo $F_1$}
{$1\ m/s^2$ hướng theo $F_2$}
\haidong{5}\loigiai{\begin{itemize}
\item ${a}=\dfrac{{F}_1-{F}_2}{\ m}=1\ m/s^2 ;\ \vec{{a}}$ hướng theo hướng của lực lớn hơn $\vec{F}_{{1}}$

\end{itemize}


} \haidong{8} \end{ex}






\subsubsection*{PHẦN 2. Thí sinh trả lời từ câu 1 đến câu 4. Trong mỗi ý a), b), c), d) ở mỗi câu, thí sinh chọn đúng hoặc sai.}
\setcounter{ex}{0}
\begin{ex}%Tailieuchuan
Một ô tô có khối lượng $360\ kg$ đang chuyển động với tốc độ $120 \ m/s$ thì người lái xe tắt máy và hãm phanh cho ô tô chuyển động thẳng chậm dần đều. Sau $17\ s$ thì tốc độ của ô tô giảm còn $63 \ km/h$.
\choiceTF[t]
{\True Gia tốc của xe ngược chiều chuyển động} 
{Gia tốc của xe có độ lớn là $0,96 \ m/s^2$} 
{Thời gian để xe dừng hẳn là $25,2 \ s$} 
{Lực hãm phanh của xe là $350\ N$} 
\loigiai{
\begin{itemchoice}
\itemch 
\itemch 
\itemch 
\itemch 
\end{itemchoice}
}
\end{ex}
\begin{ex}%Tailieuchuan
\impicinpar{Một vật có khối lượng $m=0,5 \ kg$ đang chuyển động thẳng đều trên mặt phẳng nằm ngang với vận tốc $2 \ m/s$ thì đồng thời chịu tác dụng của lực $F_1$ và $F_2=0,5 \ N$ theo phương nằm ngang như hình vẽ, vật chuyển động thẳng nhanh dần đều. Sau $4\ s$, vật đạt vận tốc $10\ m/s$.}{\begin{tikzpicture}[draw=Mapcolor, very thick,line cap=round, line join=round, scale=1, font=\footnotesize, >=stealth,draw=Mapcolor]
%[draw=Mapcolor]
\path
(0,0)coordinate(O)
;
\draw (0,0) -- (3.5,0);
\fill[pattern=north east lines] (0,0)rectangle(3.5,-0.1);
\draw  (1,0.5) rectangle (2,0);
\draw[very thick,->,Mapcolor] (2,0.25) -- (3.25,0.25);

\draw[very thick,->,Mapcolor]  (1,0.25) -- (0.25,0.25);
\node at (3.25,0.5) {$\vec F_1$};
\node at (0.25,0.5) {$\vec F_2$};
\end{tikzpicture}}
\choiceTF[t]
{\True Từ hình vẽ, ta thấy $F_1>F_2$} 
{Vật chuyển động nhanh dần đều nên gia tốc tăng dần theo thời gian} 
{\True Gia tốc của vật là $2\ m/s^2$} 
{\True Lực $F_1$ có độ lớn $1,5\ N$} 
\loigiai{\begin{itemchoice}
\itemch 
\itemch 
\itemch 
\itemch 
\end{itemchoice}
}
\end{ex}



\begin{ex}%Tailieuchuan
\immini[thm]{Một xe bán tải khối lượng $2,5$ tấn đang di chuyển trên cao tốc với tốc độ $90\ km/h$. Xe đang giữ khoảng cách an toàn so với xe ô tô chạy phía trước $70\ m$. Giả sử xe phía trước gặp sự cố và dừng đột ngột; xe bán tải tắt máy và hãm phanh, sau đó dừng lại ngay sát xe đó.}{\includegraphics[width=4cm]{images/2025_04_24_a55627a7385272a07585g-10}}
\choiceTF[t]
{\True Thực tế ta cần giữ khoảng cách an toàn lớn hơn $70\ m$ do tài xế cần một khoảng thời gian để kịp phản xạ đạp phanh}
{Gia tốc của xe bán tải xấp xỉ $5,46 \ m/s^2$} 
{\True Thời gian để xe bán tải dừng hẳn xấp xỉ $5,6\ s$} 
{Lực hãm phanh của động cơ là $12460\ N$} 
\loigiai{
\begin{itemchoice}
\itemch 
\itemch 
\itemch 
\itemch 
\end{itemchoice}
}
\end{ex}
\begin{ex}%DoanhNamDinh
\impicinpar{Một vật có khối lượng $10\ kg$ đang chuyển động thẳng đều với tốc độ $20 \ m/s$ dưới tác dụng đồng thời của hai lực $\vec{F}_1$ và $\vec{F}_2$ như hình vẽ. Độ lớn của một trong hai lực là $20\ N$.}{
\begin{tikzpicture}[line cap=round, line join=round, scale=1, font=\footnotesize, >=stealth,draw=Mapcolor]
%[draw=Mapcolor]
\path
(1.5,0.25)coordinate(O)
;


\draw[very thick,->] (1.5,0.25) -- (3.5,0.25);

\draw[very thick,->]  (1.5,0.25) -- (-0.5,0.25);
\node at (3.25,0.5) {$\vec F_1$};
\node at (0.25,0.5) {$\vec F_2$};
\draw[dashed,->] (-1.5,0.25) -- (6,0.25);
\node [Mapcolor,fill=Mapcolor,draw, thick, shape=circle, minimum width=1cm, minimum height=1cm, anchor=center] at (O) {};
\foreach \x/\g in {O/90}
\draw[fill=white](\x) circle(0.05)+(\g:0.3)node{};
\node at (6,0.75) {Hướng chuyển động};
\end{tikzpicture}}
\choiceTF[t]
{\True Độ lớn của lực còn lại là $20\ N$}
{Nếu lực $\vec{F}_1$ thôi tác dụng lên vật thì vật sẽ chuyển động thẳng đều theo hướng của lực $\vec{F}_1$}
{Nếu lực $\vec{F}_1$ thôi tác dụng lên vật thì độ lớn gia tốc của vật là $4 \ m/s^2$}
{Nếu lực $\vec{F}_1$ thôi tác dụng lên vật, quãng đường vật đi được trong $15\ s$ kể từ lúc lực $\vec{F}_1$ thôi tác dụng lên vật là $75\ m$}
\loigiai{
\begin{itemchoice}
\itemch Vật chuyển động thẳng đều, hợp lực tác dụng lên vật bằng $o \rightarrow$ hai lực ùng giá, ngược chiều, cùng độ lớn $\rightarrow F_1=F_2=20 \ N$
\itemch Nếu lực $\vec{F}_1$ thôi tác dụng lên vật thì hợp lực tác dụng lên vật là $\vec{F}_1$ khác o, vật chuyển sang trạng thái chuyển động biến đổi đều.
$\Rightarrow$ Sai
\itemch Nếu lực $\vec{F}_1$ thôi tác dụng lên vật thì độ lớn gia tốc của vật là $|a|=\left|\dfrac{F_2}{m}\right|=\dfrac{20}{10}=2 \ m/s^2$
$\Rightarrow$ Sai
\itemch Chọn chiều dương cùng chiều chuyển động, $t=o$ lúc lực $\vec{F}_1$ thôi tác dụng lên vật.
\begin{itemize}
\item gia tốc của vật $a=\dfrac{-F_2}{m}=\dfrac{-20}{10}=-2\ m/s^2$
\item thời gian từ lúc lực $\vec{F}_1$ thôi tác dụng lên vật đến lúc vật dừng lại tạm thời là
$v=v_{o}+a t=20-2 t=0 \rightarrow t=10\ s$
\item quãng đường trong 10 s đầu $S_1=v_{o} t+\dfrac{a t^2}{2}=20.10-\dfrac{2.10^2}{2}=100 \ m$
\item quãng đường trong 5 s tiếp theo (vật đổi chiều chuyển động, gia tốc không đổi)
$s_2=\left|v_{o} t+\dfrac{a t^2}{2}\right|=\left|0.5-\dfrac{2.5^2}{2}\right|=25\ m$
\item tổng quãng đường trong 15 s là $s=s_1+s_2=125\ m$
$\Rightarrow$ Sai
\end{itemize}
\end{itemchoice}
}
\end{ex}





















\subsubsection*{PHẦN 3. Thí sinh trả lời từ câu 1 đến câu 6.}
\setcounter{ex}{0} 
\begin{ex}%Tailieuchuan
Một toa tàu đang chuyển động thẳng với tốc độ không đổi. Lực cản trở tác dụng lên toa tàu có độ lớn $2000\ N$. Lực kéo toa tàu đó có độ lớn bằng bao nhiêu niutơn?
\shortans[oly]{2000}
\loigiai{}
\end{ex}
\begin{ex}%Tailieuchuan
Sau thời gian $0,02\ s$ tiếp xúc với chân của cầu thủ, quả bóng khối lượng $500\ g$ ban đầu đứng yên bay đi với tốc độ $54,0\ km/h$. Lực tác dụng lên quả bóng bằng bao nhiêu niutơn?
\shortans[oly]{375}
\loigiai{}
\end{ex}

\begin{ex}%Tailieuchuan
Một xe tải chở hàng có tổng khối lượng xe và hàng là $6$ tấn, khởi hành với gia tốc $0,4\ m/s^2$. Khi không chở hàng xe tải khởi hành với gia tốc $0,6\ m/s^2$. Biết rằng lực tác dụng vào ô tô trong hai trường hợp đều bằng nhau. Khối lượng của hàng là bao nhiêu kilogam?
\shortans[oly]{2000}
\loigiai{}
\end{ex}
\begin{ex}%book%[1P2K3-2][Loai1]
Vật đang chuyển động thẳng thì chịu tác dụng của lực $\vec F$ ngược chiều chuyển động trong $6\ s$, vận tốc giảm từ $8\ m/s$ còn $5\ m/s$. Trong $10\ s$ tiếp theo, lực tác dụng tăng gấp đôi về độ lớn còn hướng không đổi. Tốc độ của vật ở thời điểm cuối là bao nhiêu m/s? 
\shortans[oly]{5}
\haidong{5}\loigiai{
\begin{itemize}
\item Gia tốc của vật trong 6 s đầu là: $a_1 = \dfrac{5 - 8}{6} = - 0,5\ m/s^2$. 
\item Gia tốc của vật trong 10 s tiếp theo là $a_2 = 2a_1 = 2.\left(-0,5\right) = - 1\ m/s^2$.
\item Vận tốc của vật ở thời điểm cuối là: $v = 5 + a_2t = 5 - 1.10 = - 5$ m/s.
\end{itemize}
}

\haidong{8} \end{ex}
\begin{ex}%book%[1P2B3-2][Loai1]
Dưới tác dụng của lực $\vec F_1$ có độ lớn $15\ N$, vật thu được gia tốc $1\ m/s^2$. Nếu vật chịu thêm lực $\vec F_2$ có độ lớn $20\ N$ và có phương vuông góc với $\vec{F}_1$ thì gia tốc của vật có độ lớn bằng bao nhiêu $m/s^2$ (làm tròn kết quả đến chữ số hàng phần trăm)?
\shortans[oly]{1,67}
\loigiai{ 
\begin{itemize}
\item Hai lực vuông góc nên: $F'=\sqrt{F_1^2+F_2^2}=25\ N$.
\item Gia tốc tỉ lệ thuận với lực nên $\dfrac{a'}{a}=\dfrac{F'}{F}\Rightarrow a'=\dfrac{5}{3}$ N.
\end{itemize}
}	
\haidong{8} \end{ex}
\begin{ex}%Tailieuchuan
Một vật có khối lượng $0,5\ kg$ chuyển động nhanh dần đều với vận tốc ban đầu $2\ m/s$. Sau thời gian $4$ giây thì vật đi được quãng đường $24\ m$. Biết rằng vật luôn chịu tác dụng của lực kéo $F_k$
và lực cản $F_c=0,5\ N$. Nếu sau thời gian $4$ giây đó, lực kéo ngưng tác dụng thì sau bao nhiêu giây vật dừng lại?
\shortans[oly]{10}
\loigiai{}
\end{ex}





